Deputies are elected in multi-member districts (the number of seats is
published with each district's result), and voters can cast as many votes as
there are seats in their district.  To report results in a form consistent
with other countries, we normalize the number of votes cast by the number of
votes allowed per voter, counting missing votes as blank, as for Luxembourg
and Liechtenstein (`fictive voters').  Only the elections of Deputies are
included; Senators and Connétables are not.  In St Martin in 2002 the source
reports four more votes than valid ballots; they are kept as published.

Districts are drawn as their parishes or, where a district divides a parish
(St Helier, St Saviour and St Brelade), as its vingtaines, from
OpenStreetMap: St Helier No.~1 / South the Vingtaine de la Ville, No.~2 the
two Mont au Prêtre vingtaines, No.~3 \& 4 Rouge Bouillon, Mont à l'Abbé and
Mont Cochon; from 2022 North Haut du Mont au Prêtre, Mont à l'Abbé and Mont
Cochon, Central Rouge Bouillon and Bas du Mont au Prêtre.
